% ============================================================
%  Chapter 4 — The Amazing Touchscreen
%  Inside a Smartphone: The Tiny World in Your Hand
% ============================================================

\chapter{The Amazing Touchscreen}

\chapteropening{%
  Every tap, swipe, and pinch you make on a phone screen is sending
  a signal to the processor. But how does a flat piece of glass know
  where your finger is?
}
\chapterbadge{Mission: Understand the Magic of Touch!}

% --- Opening ---
\section*{Meet Tappy}

\character{Tappy}{Hello! I am Tappy — and I live right here on the surface of the screen.
  Every time a finger lands on the glass, I feel it.
  I know exactly where it touched, how many fingers touched at once,
  and how long they stayed. Then I tell Chip what to do next.}

\sectionrule

% --- How touch works ---
\section*{How Touching the Screen Sends Commands}

A smartphone screen is a \term{capacitive touchscreen}.
``Capacitive'' means it works with an invisible electrical field spread across the glass.

Your skin can conduct a tiny amount of electricity.
When your finger touches the screen:

\begin{enumerate}[label=\textbf{\arabic*.}]
  \item The screen's invisible grid of sensors detects a change in that electrical field.
  \item The sensors record the exact location of the touch.
  \item That information is sent to the processor as an instruction.
  \item The processor decides what to do (open an app, scroll the page, etc.).
\end{enumerate}

\begin{picturethis}
  Imagine the screen is covered by an invisible net of tiny tripwires.
  When your finger lands, it triggers the tripwires nearby.
  The phone measures \emph{which} wires were triggered to figure out
  exactly where you touched.
\end{picturethis}

\begin{funfact}
  Most ordinary gloves prevent touchscreens from working, because they block the electrical signal from your fingertip.
  But special ``touchscreen gloves'' are made with conductive material
  (material that lets electricity pass through)
  in the fingertips so the signal passes through.
\end{funfact}

\sectionrule

% --- Gestures ---
\section*{Swiping, Tapping, and Pinching}

Different finger movements (called \term{gestures}) send different instructions.

\begin{quickmap}
  \begin{center}
  \begin{tabular}{ll}
    \textbf{Gesture} & \textbf{What it usually does} \\
    \hline
    Tap             & Select or open something \\
    Double tap      & Zoom in or like a photo \\
    Swipe           & Scroll or switch between screens \\
    Pinch in        & Zoom out \\
    Spread out      & Zoom in \\
    Long press      & Open extra options \\
    Drag            & Move something across the screen \\
  \end{tabular}
  \end{center}
\end{quickmap}

\sectionrule

% --- Pixels and the display ---
\section*{The Screen Itself: Millions of Tiny Lights}

Beneath the touch sensors is the actual display.
It is made up of millions of tiny dots of light called \term{pixels}.
Each pixel can be set to any colour.
When millions of pixels light up together, they form the images you see.

On the same-sized screen, more pixels means a sharper, more detailed image.
The number of pixels on a screen is its \term{screen resolution}.

\begin{diagram}[A grid of pixels makes up every image on the screen.]
  \begin{tikzpicture}
    % 6 columns x 4 rows of pixels — each a hardcoded colour
    % Row 0 (bottom)
    \foreach \x/\y/\col in {
      0/0/TechPurple!85,  1/0/WarmOrange!80,  2/0/LeafGreen!75,
      3/0/CoralRed!70,    4/0/SkyBlue!80,     5/0/SoftPurple!65,
      0/1/SkyBlue!60,     1/1/TechPurple!60,  2/1/WarmOrange!60,
      3/1/LeafGreen!65,   4/1/CoralRed!55,    5/1/SunYellow!70,
      0/2/LeafGreen!50,   1/2/CoralRed!50,    2/2/SkyBlue!70,
      3/2/TechPurple!75,  4/2/WarmOrange!65,  5/2/LeafGreen!55,
      0/3/WarmOrange!70,  1/3/LeafGreen!70,   2/3/CoralRed!65,
      3/3/SkyBlue!55,     4/3/TechPurple!55,  5/3/CoralRed!40
    }{
      \fill[\col] (\x*0.8,\y*0.8) rectangle (\x*0.8+0.74,\y*0.8+0.74);
    }
    \node[below,font=\small\itshape] at (2.0,-0.25)
      {Each small square is one pixel. Together they form a picture.};
  \end{tikzpicture}
\end{diagram}

\sectionrule

\begin{newwords}
  \begin{description}[labelwidth=4.2cm, leftmargin=4.4cm]
    \item[\term{Capacitive touchscreen}] \emph{(say: kuh-PASS-ih-tiv)} — A screen that detects where your finger changes its invisible electrical field.
    \item[\term{Gesture}]               \emph{(say: JESS-cher)} — A finger movement that sends a command.
    \item[\term{Pixel}]                 \emph{(say: PIX-ull)} — A tiny dot of colour on a screen.
    \item[\term{Screen resolution}]     \emph{(say: rez-oh-LOO-shun)} — The number of pixels on a screen.
  \end{description}
\end{newwords}

\sectionrule

\begin{experiment}
  \textbf{Glove test!}

  Try tapping the screen of a smartphone with:
  \begin{enumerate}[label=\textbf{\arabic*.}]
    \item Your bare finger.
    \item A rubber glove or washing-up glove.
    \item A knitted glove.
    \item The eraser end of a pencil.
  \end{enumerate}
  Which ones work? Which ones do not? Can you explain why?

  \emph{Hint: think about whether each material conducts electricity.}
\end{experiment}

\sectionrule

\begin{bigidea}
  A touchscreen detects the \textbf{electrical charge} from your fingertip.
  Different finger gestures send different commands to the processor,
  which decides what happens on screen.
\end{bigidea}

\character{Tappy}{You are a great communicator! Now let's visit
  another amazing part of the phone — the cameras. Meet Lens in Chapter 5!}
